% !TEX root = ../main.tex
% =====================================================================
%  附录 A  流形上的定向、Spin、Spin^c、Pin^± 与 Pin^c 结构
%  附录 B  11/8 猜想
%  附录 C  关于 SW 模空间的诸定理
%  附录 D  三维流形 π1 的约束
%  原文 pp.72--78（PDF pp.73--79）
%  公式：(A.1)--(A.20)、(B.1)--(B.8)，逐条保留原文编号
%  注意：(A.9) 为原稿中的空编号，按 1:1 覆盖要求保留
% =====================================================================

这些附录收集了关于四维流形数学的补充材料。

\section{流形上的定向、自旋、$\mathrm{Spin}^{c}$、$\mathrm{Pin}^{\pm}$ 与 $\mathrm{Pin}^{c}$ 结构}

\begin{keytakeaways}
\item \textbf{物理图像：} 费米子（旋量场）只有在流形具有自旋结构时才能整体定义；自旋$^{c}$ 结构是它的最小推广，它允许费米子携带一个整体 $U(1)$ 荷——这正是 §4 中拓扑扭曲把旋量与鬼场放到四维流形上的前提。而 't Hooft 磁通所度量的，恰是 $SO(3)$ 丛不能提升为 $\mathrm{Spin}$ 丛的障碍。
\item \textbf{数学工具：} 结构群的约化（特别是沿覆盖映射的“约化”）、转移函数与 Čech 上同调、用上闭链具体构造 $w_{1}$ 与 $w_{2}$、Bockstein 同态 $W_{3}$，以及 $\mathrm{Spin}^{c}(4)\cong\{(u_{L},u_{R})\mid\det u_{L}=\det u_{R}\}$ 的表示论。
\item \textbf{核心结论：} 每个紧致可定向四维流形都带自旋$^{c}$ 结构，因为由 $\int_{\sigma}w_{2}=\sigma\cdot\sigma\bmod2$ 可知 $w_{2}$ 总有整提升；但自旋结构要求 $w_{2}=0$，而不可定向情形连 $\mathrm{Pin}^{c}$ 结构也可能不存在。
\end{keytakeaways}

\subsection{结构群的约化：一般讨论}

给定两个紧致李群 $G_{1}$、$G_{2}$ 以及同态 $\phi:G_{1}\to G_{2}$，我们可以定义一个函子 $F_{\phi}$，它把 $M$ 上的主 $G_{1}$ 丛映为 $M$ 上的主 $G_{2}$ 丛：把主 $G_{1}$ 丛 $G_{1}\to P\to M$ 映为 $(P\times_{G_{1}}G_{2})\to M$。回忆 $(P\times_{G_{1}}G_{2})$ 是二元组 $(p,g)\in P\times G_{2}$ 的集合，配以等价关系 $(ph,g)=(p,\phi(h)g)$（$h\in G_{1}$）；它显然带有一个自由的右 $G_{2}$ 作用。

\medskip
\noindent\textbf{定义}　设 $G_{2}\to P_{2}\to M$ 是一个主 $G_{2}$ 丛。在 $\phi:G_{1}\to G_{2}$ 下\emph{到 $G_{1}$ 的一个约化}，是指一个主 $G_{1}$ 丛 $G_{1}\to P_{1}\to M$ 连同同构 $\psi$，使得下图交换：
\begin{equation}
\begin{array}{ccc}
(P_{1}\times_{G_{1}}G_{2}) & \xrightarrow{\ \psi\ } & P_{2}\\[4pt]
 & \searrow & {\Big\downarrow}\\[4pt]
 & & M
\end{array}
\tag{A.1}
\end{equation}

把定义展开，可以得到用贴片重叠 $U_{\alpha\beta}$ 上的转移函数所作的描述。若 $h_{\alpha\beta}:U_{\alpha\beta}\to G_{1}$ 是 $P_{1}$ 的转移函数，则 $P_{2}$ 作为主 $G_{2}$ 丛，同构于一个以 $\phi(h_{\alpha\beta}):U_{\alpha\beta}\to G_{2}$ 为转移函数的丛。

% [Tutor Note: 一般定义与“覆盖情形”之间存在张力，读时须留意：对子群包含（下面例 1--3），“约化”给出
%  真子丛，符合直觉；但对覆盖映射（例 4），$\tilde P$ 是 $P$ 的覆盖而不是子丛，故作者明言
%  “reduction”一词在此有误导性。第 5 条把自旋结构归入后一情形——这正是“自旋结构不能看成子丛、
%  而应看成二重覆盖”的由来，也是 §4 中旋量场能整体定义的条件。]

\medskip
\noindent\textbf{例}

\begin{enumerate}[label=\arabic*., leftmargin=2.1em]
\item 若 $\phi:H\to G$ 是子群的包含，则给定 $G$-丛 $P\to M$，$H$ 自由作用，故可考虑 $G/H\to P/H\to M$，这是一个齐性空间丛。此时 $P/H$ 的一个截面给出 $P$ 到 $H$-丛的一个约化，而后者事实上是一个子丛。作为特例取 $H=\{1\}$：这就是那个熟知的事实——主 $G$ 丛的一个整体截面使该丛平凡化。

\item 取 $H=O(n)$、$G=GL(n,\RR)$。一个度量给出从标架丛到正交标架丛 $B_{O}(M)$ 的约化。显然，正交标架丛是标架丛的子丛。

\item 设 $M$ 是黎曼流形，$H=SO(n)$、$G=O(n)$，$\phi$ 为包含。则 $B_{O}(M)/H$ 是定向二重覆盖。若 $M$ 可定向，则定向丛有截面。事实上，若 $M$ 连通，$B_{O}(M)$ 有两个连通分支，因而有两个截面。选定一个截面就给出丛到定向标架的 $SO(n)$ 丛的约化；截面的选取就是 $M$ 定向的选取。

\item 现在设 $\phi$ 是覆盖映射，即 $\phi:\tilde G\to G$ 满射、核为 $K$。则主 $G$ 丛 $P$ 到主 $\tilde G$ 丛 $\tilde P$ 的“约化”是一个同构
\begin{equation}
\begin{array}{ccc}
\tilde P/K & \xrightarrow{\ \psi\ } & P\\[4pt]
 & \searrow & {\Big\downarrow}\\[4pt]
 & & M
\end{array}
\tag{A.2}
\end{equation}
注意一般定义中“约化”一词在此是有误导性的，因为 $\tilde P$ 实际上是 $P$ 的一个覆盖。换个说法：主 $G$ 丛 $P$ 到 $\tilde G$ 的“结构群约化”，等同于一个覆盖 $P$ 的主 $\tilde G$ 丛
\begin{equation}
\tilde P\to P\to M,
\tag{A.3}
\end{equation}
使得每一点 $m\in M$ 上的纤维“看起来像”覆盖
\begin{equation}
1\to K\to\tilde G\to G\to1.
\tag{A.4}
\end{equation}

\item 定向流形上自旋结构的选取是上一注记的特例，即 $\phi:\mathrm{Spin}(n)\to SO(n)$ 情形下的 (A.2)。为了分类自旋结构，我们先分类 $B_{SO}(M)$ 上的主 $\ZZ_{2}$ 丛，它们由 $z\in H^{1}(B_{SO}(M);\ZZ_{2})$ 给出。自旋结构是那些限制到纤维上给出二重覆盖 $\mathrm{Spin}(n)\to SO(n)$ 的主 $\ZZ_{2}$ 丛。二重覆盖 $\mathrm{Spin}(n)\to SO(n)$ 对应于 $z_{s}\in H^{1}(SO(n);\ZZ_{2})\cong\ZZ_{2}$。因此，自旋结构就是 $B_{SO}(M)$ 的二重覆盖中，限制到 $B_{SO}(M)\to M$ 的纤维上给出类 $z_{s}$ 的那些。注意两个自旋结构之差 $z_{1}-z_{2}$ 于是在纤维上是平凡的，从而来自 $H^{1}(M;\ZZ_{2})$ 中的一个类。因此，$M$ 上的自旋结构构成 $H^{1}(M;\ZZ_{2})$ 的一个 torsor。

\item 我们可以用这种方式处理其它结构。对流形 $M$，基于上述概念并施于同态
\begin{gather}
\phi:\mathrm{Pin}^{c}\to O(n) \tag{A.5}\\
\phi:\mathrm{Spin}^{c}\to O(n) \tag{A.6}\\
\phi:\mathrm{Pin}^{\pm}\to O(n) \tag{A.7}\\
\phi:\mathrm{Spin}\to O(n) \tag{A.8}\\
\quad \tag{A.9}
\end{gather}
就可以谈论 $\mathrm{Pin}^{c}$、$\mathrm{Spin}^{c}$、$\mathrm{Pin}^{\pm}$、$\mathrm{Spin}$ 结构。注意这些同态一般既不是单射也不是满射。对可定向流形上的 BLOTZ 结构，我们把相应的同态施于定向标架丛的 $SO(n)$。
\end{enumerate}

% [Tutor Note: 原文在 (A.8) 之后还留着一个空编号 (A.9)，其对应的公式行是空白的。这显然是一行被删去
%  而未及清理的编号（从上下文看，原本应还有一条同态）。(A.5)--(A.8) 本身恰好列举了
%  $\mathrm{Pin}^{c}$、$\mathrm{Spin}^{c}$、$\mathrm{Pin}^{\pm}$、$\mathrm{Spin}$ 四种结构，故 (A.9)
%  并无对应的数学内容。本译本按“编号 1:1 覆盖”的要求保留该空编号，绝不臆造公式。]

% [Tutor Note: “BLOTZ 结构”是作者随手写的占位名词，并非标准术语；它代表“把上述任一结构推广到
%  一般 $G$-结构”这一做法。原意是：对可定向流形，把 (A.5)--(A.8) 中的 $O(n)$ 换成 $SO(n)$，
%  再作用到定向标架丛上即可。不要把这个词当作某个已有的几何结构去查文献。]

\subsection{自旋结构与 Pin 结构的障碍}

把上述略显抽象的描述翻译成流形 $X$ 的切丛的转移函数语言，是值得做的。设 $\{U_{\alpha\beta}\}$ 是 $X$ 的坐标图册。利用度量可以构造正交标架，它们有转移函数
\begin{equation}
g_{\alpha\beta}:U_{\alpha\beta}\to O(n),
\tag{A.10}
\end{equation}
其中 $\dim_{\RR}X=n$。若我们能用一个上链把这些转移函数改造成
\begin{equation}
\tilde g_{\alpha\beta}=h_{\alpha}g_{\alpha\beta}h_{\beta}^{-1},
\tag{A.11}
\end{equation}
其中 $h_{\alpha}:U_{\alpha}\to O(n)$，使得 $\tilde g_{\alpha\beta}:U_{\alpha\beta}\to SO(n)$，则该流形可定向。可定向性的唯一障碍由一个 Čech 2-上链 $\det g_{\alpha\beta}:U_{\alpha\beta}\to\ZZ_{2}$ 给出，它定义上同调类 $w_{1}(X)\in H^{1}(X;\ZZ_{2})$。注意若 $X$ 单连通，则 $H^{1}(X;\ZZ_{2})=0$，故它必定可定向。

当 $X$ 可定向时，可取 $g_{\alpha\beta}:U_{\alpha\beta}\to SO(n)$。若选取一个\emph{好覆盖}，即所有交 $U_{\alpha\beta\cdots}$ 都是可缩的，则在每个重叠 $U_{\alpha\beta}$ 上都可以选取提升 $\tilde g_{\alpha\beta}:U_{\alpha\beta}\to\mathrm{Spin}(n)$。唯一的问题是这些提升的上链条件可能不成立。由于我们已经选了提升，而覆盖 $\mathrm{Spin}(n)\to SO(n)$ 的核恰是群 $\{\pm1\}$，我们确知在 $U_{\alpha\beta\gamma}$ 上
\begin{equation}
\tilde g_{\alpha\beta}(x)\tilde g_{\beta\gamma}(x)\tilde g_{\gamma\alpha}(x)
:=\xi_{\alpha\beta\gamma}\in\{\pm1\}\subset\mathrm{Spin}(n),
\qquad\forall x\in U_{\alpha\beta\gamma}.
\tag{A.12}
\end{equation}
符号 $\xi_{\alpha\beta\gamma}$ 定义一个 Čech 3-上链，它给出 $H^{2}(X;\ZZ_{2})$ 中的一个上同调类。这就是 $w_{2}$ 的一个（非常具体的）定义。注意：用上链对提升 $g_{\alpha\beta}$ 作任何改动、或选取不同的提升 $\tilde g_{\alpha\beta}$，都只会把 $\{\xi_{\alpha\beta\gamma}\}$ 改变一个上边缘。几乎按定义，这就是自旋结构存在的唯一障碍。

% [Tutor Note: 这里记号的偏移值得点明（否则核对文献时极易错位）：按标准 Čech 记号，转移函数定义在
%  双交叠 $U_{\alpha\beta}$ 上，是 1-上链，其上边缘落在三交叠，故 $w_{1}\in H^{1}(X;\ZZ_{2})$ 来自
%  一个 1-上闭链；作者称其为 “Čech 2-cochain” 是按指标个数（$\alpha\beta$ 两个下标）计数。
%  同理 (A.12) 的 $\xi_{\alpha\beta\gamma}$ 定义在三交叠上，是标准的 2-上闭链（作者称
%  “3-cocycle”），它给出 $w_{2}\in H^{2}(X;\ZZ_{2})$。两种记法只差一个指标的平移，数学内容不变。]

% [Tutor Note: (A.12) 是 $w_{2}$ 的“障碍论”定义的核心：$\xi_{\alpha\beta\gamma}$ 度量的是提升后的
%  转移函数在每一个三交叠上“差一个符号”的程度。它之所以自动是上闭链，是因为 $\mathrm{Spin}(n)$
%  覆盖 $SO(n)$ 的核恰为 $\{\pm1\}$ 且为中央元（与一切元素交换）；若核非中央，四交叠上会出现
%  非平凡的相容性条件。这也解释了为何 $\mathrm{Spin}^{c}$ 情形只需处理 $w_{2}$ 的整提升，
%  而不必担心更复杂的上链。]

\begin{enumerate}[label=\arabic*., leftmargin=2.1em]
\item 例：圆周上的两个自旋结构，以及与圆周的两个二重覆盖的关系。\footnote{〔原稿边注〕原文此处只写了大写提醒 “EXPLAIN.”（待解释），并未展开。}

\item 一个例子：$BSO(2)=BU(1)=\CC\PP^{\infty}$。$w_{2}$ 是 $c_{1}$ 模二的约化，而 $c_{1}$ 是与 2-胞腔对偶的上链。于是应当解释为什么 $c_{1}$ 模二是自旋结构的障碍。在 $\CC\PP^{2}$ 上 $c_{1}=3x$。复化切丛容许把结构群 $SO(4)=SU(2)\times SU(2)/\ZZ_{2}$ 约化为 $U(2)=SU(2)\times U(1)/\ZZ_{2}$。限制到 $\CC\PP^{1}\subset\CC\PP^{2}$ 上，切丛分裂为 $\mathcal{O}(2)\oplus\mathcal{O}(1)$，其中 $\mathcal{O}(2)$ 是 $\CC\PP^{1}$ 的切丛、$\mathcal{O}(1)$ 是法丛。结构群进一步约化为 $U(1)\times U(1)$。显然，与法丛 $\mathcal{O}(1)$ 相伴的主 $U(1)$ 丛不容许一个二重覆盖，使其限制到纤维上给出 $U(1)$ 到 $U(1)$ 的二重覆盖。

\item $\mathrm{Spin}^{c}$ 结构的障碍是 $W_{3}$，即 $w_{2}$ 在 Bockstein 映射下的像。因此，当 $w_{2}(X)\in H^{2}(X;\ZZ_{2})$ 有整提升时它为零。利用下述事实：对一切 $\sigma\in H_{2}(X;\ZZ)$，
\begin{equation}
\int_{\sigma}w_{2}(X)=\sigma\cdot\sigma\bmod2,
\tag{A.13}
\end{equation}
其中 $\sigma\cdot\sigma$ 是定向整相交数，可以证明这样的整提升确实存在。细节见 \cite{ref52}。类似的论断在不可定向情形下失效：$\RR\PP^{2}\times\RR\PP^{2}$ 不容许 $\mathrm{Pin}^{c}$ 结构。\footnote{〔原稿边注〕给出证明。}此外，还存在可定向的五维流形不是 $\mathrm{Spin}^{c}$ 的；一个简单例子是对称 $SU(3)$ 矩阵构成的空间，它微分同胚于 $SU(3)/SO(3)$。解释见 \cite{ref28}。
\end{enumerate}

% [Tutor Note: 〔原稿边注〕原文在此处页面边缘还有一条提醒，大意是“另可查 Gompf、Stipsicz
%  的《4-Manifolds and Kirby Calculus》中的证明”。本译本把它并入此处，以免与正文混淆。]

% [Tutor Note: $(A.13)$ 是全文最有用的一条“验算公式”：它把 $w_{2}$ 的取值归结为相交数模 2。
%  注意推论：$\int_{\sigma}w_{2}\equiv\sigma\cdot\sigma\pmod2$ 说明同态
%  $H_{2}(X;\ZZ)\to\ZZ_{2}$，$\sigma\mapsto\sigma\cdot\sigma\bmod2$，恰由 $w_{2}$ 给出；
%  这也直接给出 $w_{2}(X)$ 的整提升的存在性（因为相交形式在自由部分上幺模，取 $c_{1}$ 使
%  $c_{1}\cdot\sigma=\sigma\cdot\sigma$ 对一切整类成立即可）。四维是这一现象成立的临界维数：
%  对高维流形 $w_{2}$ 未必有整提升，故 $\mathrm{Spin}^{c}$ 结构未必存在（原文举的五维例子即此）。]

\subsection{四维流形上的 $\mathrm{Spin}^{c}$ 结构}

群 $\mathrm{Spin}^{c}(4)$ 定义为
\begin{equation}
\mathrm{Spin}^{c}(4):=\bigl(\mathrm{Spin}(4)\times U(1)\bigr)/\ZZ_{2}
=\bigl(SU(2)\times SU(2)\times U(1)\bigr)/\ZZ_{2},
\tag{A.14}
\end{equation}
其中我们按嵌入为 $(-1,-1,-1)$ 的群 $\ZZ_{2}$ 作商。$X$ 的定向正交标架丛 $OrFr(X)$ 是一个主 $SO(4)$ 丛，而自旋$^{c}$ 结构——按定义——就是沿显然的同态 $\mathrm{Spin}^{c}(4)\to SO(4)$ 把结构群约化为主 $\mathrm{Spin}^{c}(4)$ 丛。把定义展开，这意味着自旋$^{c}$ 结构由一个主 $\mathrm{Spin}^{c}(4)$ 丛 $\mathfrak{P}$ 连同投影 $\mathfrak{P}\to OrFr(X)$ 给出，而后者沿纤维看起来像正合列
\begin{equation}
1\to U(1)\to\mathrm{Spin}^{c}(4)\to SO(4)\to1.
\tag{A.15}
\end{equation}
当 $w_{2}(X)$ 有整提升时，自旋$^{c}$ 结构存在。如前所述，对每个紧致可定向四维流形这确实成立。

由 $[(v_{L},v_{R},\zeta)]\mapsto(\zeta v_{L},\zeta v_{R})$ 给出的群同态 $\mathrm{Spin}^{c}(4)\to U(2)\times U(2)$ 定义一个同构
\begin{equation}
\mathrm{Spin}^{c}(4)\cong\bigl\{(u_{L},u_{R})\;\big|\;\det u_{L}=\det u_{R}\bigr\}
\subset U(2)\times U(2);
\tag{A.16}
\end{equation}
而后一种表述使下面这一点变得显然：$\mathrm{Spin}^{c}(4)$ 有两个不等价的秩 2 复表示 $W^{\pm}$，它们直接由两个 $U(2)$ 因子的基本表示定义。于是给定一个自旋$^{c}$ 结构，就有两个相伴的复秩 2 丛 $W^{\pm}\to X$，并且如上面 §4.7.2 的讨论所述，可以等同
\begin{equation}
W^{\pm}=S^{\pm}\otimes L.
\tag{A.17}
\end{equation}
反之，这样一对丛 $W^{\pm}$ 也定义一个自旋$^{c}$ 结构。因此，自旋$^{c}$ 结构的空间是线丛群的一个 torsor：给定线丛 $\mathcal{L}$，我们总可以取
\begin{equation}
W^{\pm}\to W^{\pm}\otimes\mathcal{L},
\tag{A.18}
\end{equation}
从而
\begin{equation}
c_{1}(\det W^{\pm})\to c_{1}(\det W^{\pm})+2c_{1}(\mathcal{L}).
\tag{A.19}
\end{equation}
注意：给定 $X$ 上的一个近复结构，就存在一个标准的自旋$^{c}$ 结构
\begin{equation}
W^{+}=\Omega^{0,0}(X)\oplus\Omega^{0,2}(X),
\qquad
W^{-}=\Omega^{0,1}(X).
\tag{A.20}
\end{equation}

% [Tutor Note: (A.16) 的表示论内容值得展开。$\mathrm{Spin}(4)\cong SU(2)_{L}\times SU(2)_{R}$，故
%  $\mathrm{Spin}^{c}(4)$ 是 $U(2)\times U(2)$ 中“两个 $U(2)$ 的行列式相等”的那个子群。两个
%  $U(2)$ 因子各自的基本表示拉回，就给出 $\mathrm{Spin}^{c}(4)$ 的两个不等价复二维表示
%  $W^{+}$、$W^{-}$；它们正是旋量表示 $S^{\pm}$ 与同一个复线丛 $L$（由 $U(1)$ 因子给出）的
%  张量积，即 (A.17)。反之，给定 $\mathcal{L}$ 就同时扭动两个因子，故自旋$^{c}$ 结构在
%  “同时作 $W^{\pm}\mapsto W^{\pm}\otimes\mathcal{L}$”下等价，自旋$^{c}$ 结构的空间因此是线丛群的
%  一个 torsor。(A.19) 中的平移量是 $2c_{1}(\mathcal{L})$ 而非 $c_{1}(\mathcal{L})$——这正体现
%  “$w_{2}$ 有整提升”：$\det W^{\pm}$ 的 $c_{1}$ 模 2 只依赖自旋$^{c}$ 结构，扭动不改变它，
%  因为 $2c_{1}(\mathcal{L})$ 为偶。]

% [Tutor Note: (A.20) 给出的“标准自旋$^{c}$ 结构”是全文多次用到的“规范类”
%  $\lambda_{c}$：对近复（从而辛）流形，$W^{+}=\Omega^{0,0}\oplus\Omega^{0,2}$、
%  $W^{-}=\Omega^{0,1}$，其行列式线丛给出 $2\lambda_{c}=c_{1}(X)$。这就是 §10.1 中
%  陶布斯定理只涉及 $\pm\lambda_{c}$ 的原因（辛流形的 SW 基本类只有规范类），也是 (11.5) 中
%  “$w_{2}(X)$ 的整提升 $2\lambda_{c}$”的具体来源。]

\subsection{'t Hooft 磁通}

在讲义的若干处，我们用到主 $SO(3)$ 丛 $P$ 的第二 Stiefel--Whitney 类 $w_{2}(P)$。一般地，流形 $M$ 上的 $SO(n)$ 丛 $P$ 有一个示性类 $w_{2}(P)\in H^{2}(M;\ZZ_{2})$。定义它的一种方式沿用上面关于自旋结构的讨论：取 $M$ 的一个好覆盖 $\{U_{\alpha}\}$ 以及 $P$ 在该覆盖上的平凡化；在贴片重叠 $U_{\alpha\beta}$ 上选取提升 $\tilde g_{\alpha\beta}:U_{\alpha\beta}\to\mathrm{Spin}(n)$，并度量上链条件在 $U_{\alpha\beta\gamma}$ 上失效的程度，从而定义 $H^{2}(M;\ZZ_{2})$ 中的一个类。若 $w_{2}(P)$ 非零，则不存在 $P$ 的结构群从 $SO(n)$ 到 $\mathrm{Spin}(n)$ 的约化。

一个不提升为 $\mathrm{Spin}(n)$ 丛的 $SO(n)$ 丛的简单例子如下：取 $M$ 为二维紧致曲面，取一点 $p\in M$ 及 $p$ 周围的一个小圆盘；用绕圆盘边界的转移函数定义 $SO(n)$ 中的一个非平凡闭圈，就定义出一个 $SO(n)$ 丛。

% [Tutor Note: 这段是 §5--§7 中 't Hooft 通量的拓扑基础。$SU(2)/\ZZ_{2}=SO(3)$，故一个
%  $SO(3)$ 丛 $P$ 的 $w_{2}(P)$ 恰是“能否提升为 $SU(2)$ 丛”的障碍；在 $P$ 的提升不存在的
%  情形，规范场的第二 Stiefel--Whitney 类给出一个非平凡的“磁通”量子数，这就是 't Hooft 磁通。
%  原文的二维例子（绕边界的转移函数给出 $SO(n)$ 的非平凡闭圈）正是 $\pi_{1}(SO(n))=\ZZ_{2}$
%  的生成元在起作用：绕圆盘一周的转移函数走一个非平凡闭圈，就得到一个 $w_{2}$ 非零的丛。]

\section{11/8 猜想}

\begin{keytakeaways}
\item \textbf{物理图像：} 问题出在相交形式的“几何实现”：偶幺模形式必为 $Q=mE_{8}\oplus nH$，与 $S^{2}\times S^{2}$ 作连通和只是增加 $H$ 的个数，于是自然要问——要实现给定的 $Q$，光滑四维流形至少需要几个 $H$？11/8 猜想断言至少需要 $3|m|$ 个。
\item \textbf{数学工具：} 偶幺模格的分类（$E_{8}$ 与 $H$ 的唯一分解）、罗赫林定理（$\sigma\equiv0\bmod16$）、唐纳森的定性相交形式定理、古田的 $\mathrm{Pin}(2)$ 等变 SW 理论。
\item \textbf{核心结论：} $K3^{\#\ell}\#(S^{2}\times S^{2})^{\# t}$ 给出 $n=3\ell+t$，故 $n\ge3|m|$ 可以达到；11/8 猜想（以及针对不可约流形的更强 3/2 猜想）给出 $b_{2}$ 与 $|\sigma|$ 之间的线性界，而古田定理 $n\ge2|m|+1$ 是目前一般情形下最好的界。
\end{keytakeaways}

偶幺模形式必定是 $Q=mE_{8}\oplus nH$ 型，其中 $m$ 是任意符号的整数。注意 $E_{8}\oplus(-E_{8})\cong8H$。

罗赫林（Rokhlin）定理：若 $X$ 光滑且 $w_{2}(X)=0$，则 $\sigma(X)=0\bmod16$。等价地，若 $X$ 光滑且 $Q_{X}$ 为偶，则 $\sigma(X)=0\bmod16$。

因此，具偶不定形式的光滑流形必有 $Q=2mE_{8}\oplus nH$ 的形状。由唐纳森我们知道 $n>0$。

与 $S^{2}\times S^{2}$ 作连通和只会增加 $H$，故有趣的问题是：要使 $Q$ 成为某个\emph{光滑}四维流形 $X$ 的 $Q_{X}$，最少需要多少个 $H$？

注意
\begin{equation}
X=K3^{\#\ell}\#(S^{2}\times S^{2})^{\# t}
\qquad\Rightarrow\qquad
Q_{X}=-2\ell E_{8}\oplus(3\ell+t)H,
\tag{B.1}
\end{equation}
故此情形 $m=-\ell$、$n=3\ell+t$，其中 $\ell,t\ge0$。因此存在光滑的 $X$ 使
\begin{equation}
n\ge3|m|,
\tag{B.2}
\end{equation}
而取 $t=0$ 即达到该界。

所谓“11/8 猜想”是说，$K3^{\#\ell}$ 是对这一不等式最优的四维流形：

\medskip
\noindent\textbf{猜想}　每个具偶相交形式的光滑 $X$ 必须满足 $n\ge3|m|$；或者等价地，
\begin{equation}
b_{2}(X)\ge\frac{11}{8}|\sigma(X)|.
\tag{B.3}
\end{equation}

注意对 $Q=2mE_{8}\oplus nH$（不妨设 $n\ge0$），有
\begin{equation}
\begin{gathered}
b_{2}(Q)=16|m|+2n,\\
|\sigma(Q)|=16|m|,
\end{gathered}
\tag{B.4}
\end{equation}
故 $n\ge3|m|$ 当且仅当
\begin{equation}
b_{2}(Q)\ge22|m|=\frac{11}{8}|\sigma(Q)|.
\tag{B.5}
\end{equation}

一个更强的猜想是 3/2 猜想，它适用于\emph{不可约}的 $X$。所谓不可约流形，是指它不能分解成两个四维流形的连通和而其中一个是 $S^{4}$。

\medskip
\noindent\textbf{猜想}　设 $X$ 是具偶相交形式的光滑不可约四维流形，则
\begin{equation}
n\ge4|m|-1.
\tag{B.6}
\end{equation}
（再次注意 $m=-1$ 的 $K3$ 达到该不等式。）等价地
\begin{equation}
b_{2}(X)\ge16|m|+8|m|-2=24|m|-2
=\frac{3}{2}|\sigma(X)|-2>\frac{3}{2}|\sigma(X)|.
\tag{B.7}
\end{equation}

这方面的一些已知结果：
\begin{enumerate}[label=\arabic*., leftmargin=2.1em]
\item 定理[唐纳森]　$H$ 与 $H\oplus H$ 是 $b_{2}^{+}=1$ 或 $b_{2}^{+}=2$ 的光滑 $X$ 所能出现的仅有的偶形式。（所以若存在 $E_{8}$，则 $n\ge3$。）
\item 定理[古田]　对具偶相交形式的光滑 $X$，$n\ge2|m|+1$，即
\begin{equation}
b_{2}(X)\ge\frac{10}{8}|\sigma(X)|+2.
\tag{B.8}
\end{equation}
\end{enumerate}

古田（Furuta）等人还证明，$\pm4E_{8}\oplus5H$ 不可能是任何光滑 $X$ 的相交形式。因此，古田定理所容许、却违反 11/8 猜想的第一个未知情形是 $|m|=3$ 与 $n=7,8$。

古田的方法是利用 SW 方程的 $\mathrm{Pin}(2)$ 不变性……

% [Tutor Note: 原文此处有一处笔误须指出：作者先由罗赫林定理把光滑流形的偶不定形式写成
%  $Q=2mE_{8}\oplus nH$（这是正确的：$\sigma(2mE_{8})=16m\equiv0\bmod16$），但 (B.4) 上方又写成
%  “对 $Q=mE_{8}+nH$ 有 $b_{2}(Q)=16|m|+2n$、$|\sigma(Q)|=16|m|$”。核对可知 (B.4)、(B.5)、(B.7)
%  全部只与 $Q=2mE_{8}\oplus nH$ 自洽（因 $b_{2}(2mE_{8})=16|m|$、$|\sigma(2mE_{8})|=16|m|$），
%  故 (B.4) 中的 “$mE_{8}$” 应读作 “$2mE_{8}$”。若按字面取 $Q=mE_{8}\oplus nH$，则
%  $b_{2}=8|m|+2n$、$|\sigma|=8|m|$，(B.5) 与 (B.7) 都不再成立。]

% [Tutor Note: 两条猜想的数值可以随手验算，以确认没有抄错。11/8 猜想：
%  $b_{2}=16|m|+2n\ge16|m|+6|m|=22|m|=(11/8)\cdot16|m|=(11/8)|\sigma|$ ✓。
%  3/2 猜想：$n\ge4|m|-1$ 给出 $b_{2}\ge16|m|+8|m|-2=24|m|-2=(3/2)|\sigma|-2$ ✓。
%  古田定理的界 $b_{2}\ge(10/8)|\sigma|+2$ 在 $Q=\pm4E_{8}\oplus5H$（即 $m=2$、$n=5$、$b_{2}=42$、
%  $|\sigma|=32$）处恰好取等：$(10/8)\cdot32+2=42$ ✓，而 11/8 猜想要求 $n\ge6$ 即 $b_{2}\ge44$，
%  故该形式违反 11/8 而未被古田定理排除——这正是原文说它是“第一个未知情形”之前的最后一个例子。
%  $K3$（$m=-1$，$n=3$）则同时达到两条猜测的界。]

% [Tutor Note: 末句 “The method of 古田 is to exploit a Pin(2) invariance of the SW equations....”
%  在原文中就以省略号终止，是一条未写完的句子（原稿显然还想写古田不等式的证明思路）。
%  本译本照原样以省略号结束，不代为补全。就数学内容而言，古田的做法是对 SW 方程施加
%  $\mathrm{Pin}(2)$ 等变结构，用等变 K-理论/等变上同调得到比 11/8 更弱的界 (B.8)；
%  这一证明的技术细节不在本讲义范围内。]

\section{关于 SW 模空间的诸定理}

\begin{keytakeaways}
\item \textbf{物理图像：} SW 模空间是红外希格斯分支的载体：它是 SW 方程的解模去规范群后的空间，其维数由指标（鬼数）决定，而其紧致性、有限性与定向则由椭圆算子的弗雷德霍姆性质与 Weitzenböck 型恒等式保证。
\item \textbf{数学工具：} 弗雷德霍姆算子与指标、横向性（generic metric / generic path 上的光滑性）、模空间之间的配边、以及由 $H^{1}\oplus H^{2,+}$ 上的定向诱导的模空间定向。
\item \textbf{核心结论：} 八条标准定理；其中单型猜想（第 6 条）断言单连通且 $b_{2}^{+}\ge2$ 时模空间只能是零维的，第 7 条给出辛流形上的证明，第 8 条则用反例说明这些前提条件缺一不可。
\end{keytakeaways}

\begin{enumerate}[label=\arabic*., leftmargin=2.1em]
\item $b_{2}^{+}\ge1$、一般度量：$\mathcal{M}(\lambda)$ 或者为空，或者是预期维数的光滑流形。\footnote{〔原稿边注〕原文于本节开头大写提醒：“CITE REFERENCES! AT LEAST write the key Weitzenbock identity”（补引文献！至少写出关键的 Weitzenböck 恒等式）。}

\item $b_{2}^{+}\ge2$：对一般的路径 $g_{t}$，我们得到光滑流形之间的光滑配边。

\item $\mathcal{M}(\lambda)$ 只对有限多个自旋$^{c}$ 结构非空。

\item $\mathcal{M}(\lambda)$ 若非空则紧致。

\item $\mathcal{M}(\lambda)$ 可定向，其定向依赖于 $H^{1}\oplus H^{2,+}$ 上的一个定向（与唐纳森理论中相同）。

\item \textbf{单型猜想：} 对任何单连通且 $b_{2}^{+}\ge2$ 的 $X$，若 $\mathcal{M}(\lambda)$ 非空，则它是零维的。

\item 单型猜想对辛流形成立。

\item 存在 $b_{2}^{+}=1$ 以及非单连通 $X$ 的例子，说明 $\mathcal{M}(\lambda)$ 可以有任意高的维数。
\end{enumerate}

% [Tutor Note: 这些定理的“物理来源”值得点明。模空间的预期维数由指标给出：
%  $\dim\mathcal{M}(\lambda)=n(\lambda)=\bigl(\lambda^{2}-(2\chi+3\sigma)\bigr)/4$（即 §5 中的鬼数），
%  故它只依赖 $\lambda$ 与流形的拓扑；而第 1、2、4 条（光滑性、配边、紧致性）都来自同一套
%  Fredholm 理论：SW 算子的线性化是椭圆算子，其指标有限、核的维数由横向性控制；
%  紧致性则来自 Weitzenböck 型恒等式——把方程平方求和后，$|\nabla\psi|^{2}$ 与曲率项、标量曲率
%  的项相联系，从而给出解的一致界（原文要求“至少写出”的正是这个恒等式）。]

% [Tutor Note: 第 5 条的定向须仔细：唐纳森理论中模空间的定向只依赖 $H^{2,+}(X)$ 的定向，而
%  SW 情形还要 $H^{1}(X)$ 的定向。原因是 SW 模空间是“旋量丛的截面零点 + 联络模空间”的纤维积，
%  狄拉克算子的行列式线丛的定向由 $H^{1}(X)$（即 $\ker D_{A}$ 的指标零模）决定；对单连通
%  $X$，$H^{1}=0$，于是回到熟悉的唐纳森情形。这一点在第 6 条“单连通”假设下尤为重要。]

% [Tutor Note: 第 6--8 条之间的逻辑关系是本节的关键。单型猜想（第 6 条）若成立，则 SW 不变量
%  就是有限个带符号点的计数（“字面意义地数方程的解”），这正是 §9 中“SW 单型”的定义；
%  第 7 条（辛流形）由陶布斯的 SW=GW 定理给出（见 §10.1），因为伪全纯曲线的计数给出
%  零维模空间；第 8 条说明：一旦去掉“单连通”或允许 $b_{2}^{+}=1$，模空间可以有任意高的维数，
%  故第 6 条的两个假设都不是多余的。]

\section{三维流形 $\pi_{1}$ 的约束}

\begin{keytakeaways}
\item \textbf{物理图像：} 三维流形正是 §5 中 $X=\Sigma_{3}\times\RR$ 的类空切片；其基本群受黑格分解的强约束——存在一个生成元与关系一样多的展示。若万有覆盖还是有理同调球，则 $\pi_{1}$ 必须能自由作用在 $S^{3}$ 上，于是其分类群的上同调被迫具有周期性。
\item \textbf{数学工具：} 黑格分解、塞弗特--范坎彭定理、万有覆盖与 Borel 构造 $E G\times_{\pi_{1}}S^{3}$、球丛的 Gysin 序列、Thom 同构、以及群上同调的四重周期性。
\item \textbf{核心结论：} Gysin 序列迫使 $H^{*}(B\pi_{1})$ 在足够高的次数上四重周期；而 $(\ZZ_{2})^{n}$ 的上同调并不周期，故 $(\ZZ_{2})^{n}$ 不能是任何三维流形的 $\pi_{1}$。这一论证是 Graeme 西格尔向作者解释的。
\end{keytakeaways}

任意可定向三维流形都容许一个黑格（Heegaard，原文拼作 Heegard）分解。这意味着我们可以找到一个嵌入的闭可定向曲面 $\Sigma\subset Y$，使得 $Y$ 是把两个把手体（即 $\Sigma$ 到空集的配边）用一个 $\Sigma$ 的微分同胚粘合起来的结果。这样的黑格分解强有力地约束了基本群，依据是塞弗特--范坎彭定理：一组生成元可由两个把手体各自的生成元给出。亏格为 $g$ 的 $\Sigma$ 所对应的把手体，其基本群恰是 $g$ 个生成元上的自由群。（“$A$-圈”收缩为一点，只留下“$B$-圈”而没有任何关系。）粘合到另一个把手体之后，会出现 $g$ 个关系，表达第二个把手体的新“$A$-圈”的可缩性。

于是，对任何可定向三维流形 $Y$，$\pi_{1}(Y,y_{0})$ 容许一个展示，其中生成元的个数等于关系的个数。这样的展示对一般的有限生成群显然是不可能的。

若 $Y$ 的万有覆盖 $\tilde Y$ 是有理同调球，我们还能多说一些。\footnote{这一点是 Graeme 西格尔向我解释的。}（由庞加莱猜想我们知道它就是 $S^{3}$，但下面的论证实际上并不需要这一点。）于是 $\pi_{1}$ 必须连续作用在 $S^{3}$ 上。

现在考虑 $EG\times_{\pi_{1}}S^{3}$ 的 Gysin 序列。一方面，它同伦等价于 $S^{3}/\pi_{1}$，而后者的上同调不可能在次数 3 以上非零。另一方面，它是 $BG$ 上的一个球丛。再考虑对 $(D,S)$ 的 Gysin 序列，其中 $S$ 是 $BG$ 上上述球丛，$D$ 是相应的圆盘丛。（把 $\pi_{1}$ 的作用延拓到圆盘上。）

当然，圆盘丛可缩到 $BG$，而由 Thom 同构 $H^{*}(D,S)\cong H^{*+4}(BG)$；于是在 Gysin 序列中得到 $H^{k}(BG)\to H^{k+4}(BG)$，而它被零夹在中间（对 $k>3$，因为 $H^{k}(S)\cong H^{k}(S^{3}/\pi_{1})$）。事实上，$H^{k}(BG)\to H^{k+4}(BG)$ 是乘以欧拉类。这说明：在足够高的次数以上，$BG$ 的上同调必须是周期的。但对足够多个 $\ZZ_{2}$ 的乘积，我们已经看到其上同调不是四重周期的。因此，这样的群不可能作为任何三维流形的 $\pi_{1}$ 出现。

% [Tutor Note: 这段 Gysin 序列论证值得逐步写清，因为原文压缩得极紧。记 $G=\pi_{1}$、
%  $B=BG=K(G,1)$。Borel 构造 $E G\times_{G}S^{3}$ 是 $B$ 上的球丛（纤维 $S^{3}$），且因 $EG$
%  可缩，它同伦等价于 $S^{3}/G=Y$，故其 $3$ 次以上上同调为零。取相应的圆盘丛 $D$（把作用
%  延拓到圆盘上），则 $(D,S)$ 给出 Thom 同构 $H^{k}(D,S)\cong H^{k+4}(B)$，而 $D$ 可缩到 $B$。
%  Gysin 正合列
%  \[
%   \cdots\to H^{k}(B)\xrightarrow{\ \cup e\ }H^{k+4}(B)\to H^{k+4}(S)
%   \to H^{k+1}(B)\xrightarrow{\ \cup e\ }H^{k+5}(B)\to\cdots
%  \]
%  中，当 $k$ 足够大时 $H^{k+4}(S)=0$（因 $H^{*}(S)\cong H^{*}(S^{3}/G)$ 在次数 $>3$ 为零），
%  故 $\cup e:H^{k}(B)\to H^{k+4}(B)$ 满射；同时 $H^{k+1}(S)=0$ 又给出它是单射——两端零夹住的
%  箭头必是同构。于是 $H^{k}(B)\cong H^{k+4}(B)$ 对一切足够大的 $k$ 成立，即 $H^{*}(B;R)$
%  四重周期（$R$ 取 $\ZZ$ 或 $\mathbb{F}_{2}$ 皆可）。另一方面，取 $G=(\ZZ_{2})^{n}$（$n\ge2$），
%  则 $H^{*}(B(\ZZ_{2})^{n};\mathbb{F}_{2})=\mathbb{F}_{2}[x_{1},\dots,x_{n}]$（$|x_{i}|=1$），
%  $\dim H^{k}=\binom{n+k-1}{k}$ 随 $k$ 多项式增长，绝不周期；矛盾。故 $(\ZZ_{2})^{n}$ 不能是
%  三维流形的 $\pi_{1}$。]

% [Tutor Note: 反过来看，这条论证的“逆”正是经典定理：一个有限群若能自由作用在球面上（从而
%  上同调周期），则它是球面自由作用群——Swan 的周期上同调理论给出，四重周期上同调的有限群
%  恰好是 $SU(2)$ 的有限子群（循环群、双面体群、二元正四面体群、二元八面体群、二元二十面体群）。
%  附录 D 的论证可以看作这一代数定理的纯几何版本。也正因如此，“$A$-圈的数目等于 $B$-圈的数目”
%  这一展示性质（亏格 $g$ 的两个把手体给出 $g$ 个关系）远不足以刻画三维流形的基本群。]

% [Tutor Note: 关于“万有覆盖是有理同调球”这一假设：$\tilde Y$ 按定义单连通，若它是有理同调球，
%  则由庞加莱猜想 $\tilde Y\cong S^{3}$。但整段论证实际只用到 $H^{k}(S^{3}/G)=0$（$k>3$，
%  有理系数即可），故作者的括号注“下面的论证实际上不需要这一点”是对的：只需
%  $Y=S^{3}/G$ 在次数 $3$ 以上上同调为零。原文“we can same more” 是 “say more” 的笔误。]
